% 原书第77页：只替换本页题图，保留全部已有补充方法。
\QHeading{第1题　按要求画一画。}
\PSeventySevenFirst{1}{1}
\par\vspace{3mm}
\QHeading{第1题①　补成轴对称图形的另一种作法}
\Answer{除逐点数格外，可以沿题目给定的对称轴折叠描画。两法画出的对应点相同：对应点的连线垂直于对称轴，且到轴的距离相等。}
\QHeading{第1题②、第2、3题　运动作图的不同办法}
\Answer{旋转：可用量角器逐点转过90度并保持到旋转中心的距离，也可在方格纸上按直角转向数格。平移：可先移动所有顶点再连接，也可用透明纸保持图形方向不变整体平移。先按题目要求的方向和格数完成，另一方法用来核对；移动顺序不同不另算一个最终答案。}
\QHeading{第2题　填一填，画一画。}
\noindent\makebox[\linewidth][c]{\PSeventySevenFish{1}}\par\vspace{2mm}
\PlainItem{① 小鱼从位置$A$游到位置$B$，需要先向右平移（\textcolor{answerred}{8}）格，再向（\textcolor{answerred}{上}）平移（\textcolor{answerred}{3}）格。}
\PlainItem{② 画出小鱼从位置$A$先向下平移3格，再向右平移8格后得到的图形。}
% 与修订前一致，第3题仍在该原书页的第2张答案页，不改变后续页码。
\clearpage
\QHeading{第3题　画出把下面的图形向右平移11格后的图形。}
\noindent\makebox[\linewidth][c]{\PSeventySevenTree{1}}\par
